\begin{oving}
\begin{itemize}[leftmargin=5.6em,labelwidth=5em,labelsep=0.6em,align=left]
\item[\lett] \href{https://github.com/gustavbv/Matte4C---eksamen/blob/main/oppgavesett/F10-varmelikning-og-differanser/2-Bergmann-eksamener/016-TMA4125-2024-var-oppg5-oppgave.pdf}{Matte 4N vår 2024, oppg.~5} {\color{graa}(\href{https://github.com/gustavbv/Matte4C---eksamen/blob/main/oppgavesett/F10-varmelikning-og-differanser/2-Bergmann-eksamener/016-TMA4125-2024-var-oppg5-svar.pdf}{LF})}: d'Alemberts formel på $\R$
\item[\lett] \href{https://github.com/gustavbv/Matte4C---eksamen/blob/main/oppgavesett/F10-varmelikning-og-differanser/1-TMA4432-ovinger-og-proveeksamen/003-TMA4432-oving-10-oppg1-oppgave.pdf}{Øving 10, oppg.~1} {\color{graa}(\href{https://github.com/gustavbv/Matte4C---eksamen/blob/main/oppgavesett/F10-varmelikning-og-differanser/1-TMA4432-ovinger-og-proveeksamen/003-TMA4432-oving-10-oppg1-svar.pdf}{LF})}: d'Alembert med $f=\sin x$, $g=\ee^{-x/4}$
\item[\middels] \href{https://github.com/gustavbv/Matte4C---eksamen/blob/main/oppgavesett/F10-varmelikning-og-differanser/3-hele-oppgaven-i-pensum/033-TMA4130-35-2024-host-oppg6-oppgave.pdf}{TMA4130 høst 2024, oppg.~6} {\color{graa}(\href{https://github.com/gustavbv/Matte4C---eksamen/blob/main/oppgavesett/F10-varmelikning-og-differanser/3-hele-oppgaven-i-pensum/033-TMA4130-35-2024-host-oppg6-svar.pdf}{LF})}: separert løsning sammenlignet med d'Alembert
\item[\vanskelig] \href{https://github.com/gustavbv/Matte4C---eksamen/blob/main/oppgavesett/F10-varmelikning-og-differanser/3-hele-oppgaven-i-pensum/031-TMA4125-30-35-2024-kont-oppg7-oppgave.pdf}{Matte 4N kont.\ 2024, oppg.~7} {\color{graa}(\href{https://github.com/gustavbv/Matte4C---eksamen/blob/main/oppgavesett/F10-varmelikning-og-differanser/3-hele-oppgaven-i-pensum/031-TMA4125-30-35-2024-kont-oppg7-svar.pdf}{LF})}: d'Alembert via Fouriertransformen
\end{itemize}
\end{oving}
